\section{Descomposición de Benders}
\label{sec:benders}

Esta sección corresponde a la \textbf{Entrega 3}: maestro, subproblema primal y dual, derivación
de los cortes y justificación de validez.

\subsection{Recordatorio: Benders genérico (forma del curso) y su especialización al pMP}
El curso (Cap.~3) plantea Benders sobre la forma canónica
$\min_{x\in X,\,y\ge0}\ c^\top x + d^\top y$ s.a. $Ax+By\ge b$, con $x$ \emph{complicantes}
(discretas) e $y$ operacionales. Fijado $\bar x$, el subproblema en $y$ es un LP; su función de
recurso $\theta(x)=\min\{d^\top y: By\ge b-Ax,\ y\ge0\}$ es \emph{convexa} (máximo de funciones
lineales en su dual), y el maestro la reemplaza por una variable $\eta\ge\theta(x)$:
\[
  \min_{x\in X,\,\eta}\ c^\top x+\eta,\qquad \eta\ge\theta(x).
\]
El subproblema dual vive en el poliedro $\Pi=\{\pi\ge0:B^\top\pi\le d\}$, que \emph{no depende de
$\bar x$}. Por \textbf{dualidad fuerte}, si el subproblema es factible y acotado, un \textbf{punto
extremo} $\pi^\star$ de $\Pi$ da el \textbf{corte de optimalidad} $\eta\ge\pi^{\star\top}(b-Ax)$;
si el subproblema es infactible, un \textbf{rayo extremo} $\hat\pi$ da el \textbf{corte de
factibilidad} $\hat\pi^\top(b-Ax)\le0$. Como $\Pi$ tiene un número \emph{finito} de puntos y
rayos extremos, el algoritmo termina en finitas iteraciones.

\paragraph{Especialización al pMP.} Aquí las complicantes son las aperturas $y$ (el ``$x$'' del
curso) y el recurso se \emph{desagrega por cliente}: $\eta\equiv\sum_{i=1}^N\theta_i$, donde cada
$\theta_i$ aproxima la distancia de asignación del cliente $i$. El resto de la sección instancia
esta plantilla y muestra que, por la estructura del pMP, el subproblema es siempre factible y
acotado (solo puntos extremos $\Rightarrow$ solo cortes de optimalidad).

\subsection{Qué se fija, qué se elimina, qué aparece}
Partimos de \textbf{F3}. Fijado $y=\bar y$, F3 se separa en $N$ subproblemas independientes (uno
por cliente), cada uno calculando la distancia de asignación de ese cliente. Las variables $z$
son la parte ``fácil''. Por lo tanto:
\begin{itemize}[noitemsep]
  \item Se \textbf{eliminan del maestro} todas las variables $z^k_i$.
  \item Se \textbf{introduce una variable continua} $\theta_i\ge 0$ por cliente, que representa su
        distancia de asignación.
\end{itemize}

\subsection{Problema maestro (MP)}
\begin{equation}
  \min \sum_{i=1}^N \theta_i
  \quad\text{s.a.}\quad
  \sum_{j=1}^M y_j=p,\quad
  \theta_i\ \text{satisface } BD_i\ (i\in[N]),\quad
  y_j\in\{0,1\},
  \label{eq:master}
\end{equation}
donde $BD_i$ es el conjunto de cortes de Benders del cliente $i$, inicialmente vacío y que crece
con las iteraciones.

\begin{intuicion}
El maestro es pequeño: solo decide qué sitios abrir ($y$) y ``estima'' la distancia de cada
cliente con $\theta_i$. Los cortes van obligando a que $\theta_i$ no subestime la distancia real.
\end{intuicion}

\subsection{Subproblema primal \texorpdfstring{$SP_i(\bar y)$}{SPi}}
Para $\bar y$ fijo, el subproblema del cliente $i$ (derivado de F3) es
\begin{align}
  \min\ & D^1_i + \sum_{k=1}^{K_i-1}\big(D^{k+1}_i-D^k_i\big)z^k_i \nonumber\\
  \text{s.a.}\ & z^1_i \ge 1-\!\!\sum_{j:\,d_{ij}=D^1_i}\!\!\bar y_j \label{eq:sp1}\\
  & z^k_i - z^{k-1}_i \ge -\!\!\sum_{j:\,d_{ij}=D^k_i}\!\!\bar y_j && k\in\{2,\dots,K_i\} \label{eq:spk}\\
  & z^k_i\ge 0 && k\in[K_i]. \nonumber
\end{align}
Es \textbf{siempre factible} (basta tomar $z$ suficientemente grande) y \textbf{acotado} (objetivo
con coeficientes $\ge 0$ y minimización). Esta propiedad es la que determina el tipo de corte.

\subsection{Subproblema dual \texorpdfstring{$DSP_i(\bar y)$}{DSPi}}
Con variables duales $v^k_i\ge0$ (una por restricción primal):
\begin{align}
  \max\ & D^1_i + v^1_i\Big(1-\!\!\sum_{j:\,d_{ij}=D^1_i}\!\!\bar y_j\Big) - \sum_{k=2}^{K_i} v^k_i\!\!\sum_{j:\,d_{ij}=D^k_i}\!\!\bar y_j \nonumber\\
  \text{s.a.}\ & v^k_i - v^{k+1}_i \le D^{k+1}_i-D^k_i && k\in[K_i-1] \label{eq:dspcon}\\
  & v^k_i\ge 0 && k\in[K_i]. \nonumber
\end{align}

\subsection{Derivación del corte de Benders}
De un punto extremo $\bar v$ del dual se obtiene el \textbf{corte de optimalidad}, lineal en $y$:
\begin{equation}
  \theta_i \ge D^1_i + \bar v^1_i\Big(1-\!\!\sum_{j:\,d_{ij}=D^1_i}\!\! y_j\Big)
  - \sum_{k=2}^{K_i}\bar v^k_i\!\!\sum_{j:\,d_{ij}=D^k_i}\!\! y_j.
  \label{eq:cut16}
\end{equation}

\textbf{Qué significa.} \eqref{eq:cut16} es una \emph{cota inferior lineal de} $\theta_i$ \emph{en
función de} $y$: ``tu estimación de la distancia del cliente $i$ no puede ser menor que esta
expresión, que depende de qué sitios abras''. A medida que se agregan cortes para distintos
$\bar y$, el maestro obliga a $\theta_i$ a igualar la distancia real.

\subsection{Solo cortes de optimalidad: justificación de validez}
\label{sub:optvsfeas}
En la teoría general de Benders hay dos tipos de corte:
\begin{itemize}[noitemsep]
  \item \textbf{Cortes de factibilidad:} eliminan un $\bar y$ que vuelve \emph{infactible} el
        subproblema (provienen de rayos extremos del dual no acotado).
  \item \textbf{Cortes de optimalidad:} corrigen una \emph{subestimación} del costo (provienen de
        puntos extremos del dual acotado).
\end{itemize}

\begin{quote}
\textbf{Resultado clave.} Como $SP_i(\bar y)$ es \emph{siempre factible y acotado} para todo
$\bar y$ (Sección anterior), su dual $DSP_i(\bar y)$ es \emph{siempre acotado}: nunca tiene rayos
extremos. Por lo tanto \textbf{nunca se requieren cortes de factibilidad}; todos los cortes
\eqref{eq:cut16} son de optimalidad.
\end{quote}

Esto simplifica notablemente la implementación: el separador solo debe verificar si
$\bar\theta_i$ subestima $\OPT(SP_i(\bar y))$ y, de ser así, agregar \emph{un} corte de
optimalidad. La \textbf{validez global} del método se sigue de que (i) cada corte es una
desigualdad válida (cota inferior correcta de $\theta_i$), de modo que el maestro nunca excluye la
solución óptima; y (ii) existe un número \emph{finito} de cortes distintos de la forma
\eqref{eq:cut16} (uno por cada $\bar y$ binario relevante), por lo que el branch-and-cut con
cortes perezosos termina con el óptimo exacto.

\medskip
La forma cerrada de $\bar v$ (Sección~\ref{sec:separacion}) convierte \eqref{eq:cut16} en una
desigualdad que se construye sin resolver el dual explícitamente. Notablemente, las
desigualdades resultantes coinciden con las que se obtienen sobre F1 (Cornuejols et al., 1980;
Magnanti \& Wong, 1981), aunque aquí se derivaron desde F3.

\subsection{¿Por qué este Benders es mejor que el Benders clásico?}
\label{sub:porque-mejor}
El Benders clásico (Cap.~3) resuelve, en cada iteración, un \emph{LP} de subproblema para obtener
$\pi^\star$ y el corte. La propuesta del paper mejora esto en cinco aspectos concretos:
\begin{enumerate}[leftmargin=1.6em]
  \item \textbf{Parte de una formulación fuerte.} La descomposición se construye sobre \textbf{F3},
        cuya relajación LP es ajustada y cuya matriz es rala; el maestro hereda esa calidad.
  \item \textbf{El subproblema tiene estructura especial.} No es un LP genérico: es el cálculo de
        la distancia de un cliente, con una cadena de restricciones $z^k_i-z^{k-1}_i\ge\cdots$
        cuya solución óptima es monótona.
  \item \textbf{No se resuelve ningún LP.} Esa estructura permite obtener el dual óptimo
        $\bar v$ por fórmula (complementariedad), evitando llamar al solver por cada corte.
  \item \textbf{Separación en forma cerrada $O(NM)$.} Calcular $\tilde k_i$ por cliente cuesta
        $O(M)$; sobre $N$ clientes, $O(NM)$ por separación
        (Sección~\ref{sec:separacion}). El Benders clásico no tiene esta complejidad lineal.
  \item \textbf{Solo cortes de optimalidad.} Como $SP_i$ es siempre factible y acotado, no hay
        rayos extremos ni cortes de factibilidad: la lógica de separación es más simple y robusta.
\end{enumerate}
Además, los cortes no se generan en un bucle externo maestro--subproblema, sino \emph{dentro} del
árbol de branch-and-bound mediante un callback de cortes perezosos: esto es
\textbf{branch-and-Benders-cut} (Cap.~3), que evita reoptimizar el MIP desde cero en cada
iteración.
